\begin{table}[h!]
\centering
\begin{tabular}{| p{0.20\textwidth} | p{0.50\textwidth} | p{0.30\textwidth} |}
\hline
\textbf{Tahap} 
& \textbf{Deskripsi} 
& \textbf{Tujuan} \\
\hline

Verifikasi Implementasi
& Menguji modul \textit{preprocessing}, TF-IDF, \textit{encoding}, \textit{training}, dan evaluasi pada subset kecil data. Menganalisis sampel jalur input-output secara manual.
& Memastikan \textit{pipeline} berjalan lancar dan sesuai desain. \\
\hline

Pembagian Data
& \textit{Train-test split} atau \textit{k-fold cross validation} dengan proporsi kelas yang seimbang.
& Membentuk basis evaluasi yang valid dan terkontrol. \\
\hline

Eksperimen per Model
& Melatih NB, SVM, dan RF pada skenario S1 dan S2. Mencatat metrik kinerja dan \textit{confusion matrix}.
& Membandingkan model pada skenario fitur yang berbeda. \\
\hline

Evaluasi Performa
& Menghitung \textit{accuracy}, \textit{precision}, \textit{recall}, dan \textit{F1-score}. Melakukan analisis \textit{trade-off} kesalahan model.
& Menilai kualitas prediksi secara komprehensif. \\
\hline

Penentuan Model Terbaik
& Menetapkan \textit{F1-score} sebagai metrik utama. \textit{Precision} dan \textit{recall} dianalisis untuk memahami risiko kesalahan.
& Mengidentifikasi model yang paling efektif. \\
\hline

Validasi Hasil
& Mengulangi eksperimen dengan \textit{seed} berbeda atau membandingkan pola hasil dengan literatur.
& Menjamin hasil stabil dan dapat dipercaya. \\
\hline

\end{tabular}
\caption{Tahapan Pengujian dan Evaluasi \textit{pipeline} Deteksi Hoaks}
\label{tbl:pengujian-evaluasi}
\end{table}
